
\subsubsection{Overview}

ProvenanceCache is a retrieval-aware KV cache eviction policy for long-context RAG systems. Given a query $q$ and a set of retrieved passages $\{p_1, \ldots, p_n\}$ with relevance scores $\{s_1, \ldots, s_n\}$ from a dense retriever (Contriever), ProvenanceCache allocates cache budget across three tiers based on retrieval provenance, then applies diversity-aware selection within each tier.

\subsubsection{Provenance-Aware Tiered Eviction}

The concatenated context is partitioned into three tiers based on retrieval metadata:

\textbf{Tier 0 (Query Tokens)}: All tokens from the original query $q$ receive highest retention priority. Query tokens serve as attention anchors during generation. Tier 0 receives 10\% of total cache budget $B$.

\textbf{Tier 1 (High-Relevance Passages)}: Passages with retrieval scores above median: $P_{\text{high}} = \{p_i : s_i \geq \text{median}(s_1, \ldots, s_n)\}$. These passages contain semantically relevant evidence identified by the retriever. Tier 1 receives 60\% of budget $B$.

\textbf{Tier 2 (Low-Relevance Passages)}: Passages with below-median scores: $P_{\text{low}} = \{p_i : s_i < \text{median}(s_1, \ldots, s_n)\}$. These passages preserve contrastive evidence and negative examples needed for multi-hop reasoning. Tier 2 receives 30\% of budget $B$.

The tier allocation (10/60/30) was validated empirically via grid search over {5\%, 10\%, 15\%} for Tier 0 and {50\%, 60\%, 70\%} for Tier 1, with Tier 2 absorbing the remainder.

\subsubsection{Diversity-Aware Selection via MMR}

Within each passage tier (Tiers 1 and 2), Maximal Marginal Relevance (MMR) selection is applied to prevent redundant passage retention. Given a passage set $P$, current selected set $S$, and remaining budget $k$, we iteratively select:

$$
p^* = \arg\max_{p \in P \setminus S} \left[ \lambda \cdot s_p - (1-\lambda) \cdot \max_{p' \in S} \text{sim}(p, p') \right]
$$

where:
\begin{itemize}
\item $s_p$ is the Contriever relevance score (dot-product similarity between query and passage embeddings)
\item $\text{sim}(p, p')$ is the cosine similarity between passage embeddings (Contriever 768-dim vectors)
\item $\lambda = 0.5$ balances relevance and diversity
\end{itemize}

The diversity term $\max_{p' \in S} \text{sim}(p, p')$ penalizes passages semantically similar to already-selected passages, spreading cache budget across complementary evidence sources.

\subsubsection{Integration with Llama-2-7B}

ProvenanceCache integrates with Llama-2-7B via the following pipeline:

\textbf{1. Preprocessing}: After retrieving passages via Contriever, extract:
\begin{itemize}
\item Passage boundaries (start/end token positions in concatenated context)
\item Contriever relevance scores (query-passage dot products)
\item Contriever passage embeddings (for diversity computation)
\end{itemize}

\textbf{2. Tier Assignment}: Map each token position to its tier (0/1/2) based on whether the token belongs to the query, a high-relevance passage, or a low-relevance passage.

\textbf{3. MMR Selection}: Within Tiers 1 and 2, apply MMR to select diverse passage subsets that fit within tier budgets. Evict tokens from non-selected passages.

\textbf{4. Prefill + Generation}: Run standard Llama-2-7B inference with the reduced KV cache. The eviction is applied before the first generation token.

\textbf{Computational Overhead}: MMR diversity computation requires $O(n^2)$ passage similarity computations, where $n$ is the number of retrieved passages (typically 5-10). For $n=10$, this adds ~5-10ms overhead per query on CPU. Passage embeddings are cached from retrieval preprocessing.

\subsubsection{Baseline: H2O Heavy-Hitter Oracle}

H2O \citep{zhang2023h2o} is used as the baseline. H2O tracks accumulated attention scores across all layers during prefill:

$$
\text{score}_i = \sum_{\ell=1}^{L} \sum_{j=1}^{N} \text{Attn}_\ell(i, j)
$$

where $\text{Attn}_\ell(i, j)$ is the attention weight from position $j$ to position $i$ at layer $\ell$. After prefill, H2O retains the top-$k$ tokens by accumulated score plus a sliding window of recent tokens.

H2O requires a full prefill pass to compute attention scores before eviction. ProvenanceCache evicts based on retrieval metadata before generation, avoiding redundant computation for passages destined for eviction.

\subsubsection{Theoretical Motivation}

ProvenanceCache rests on two hypotheses validated in experiments:

\textbf{H1 (Retrieval-Attention Alignment)}: Dense retrieval scores correlate moderately (Spearman $\rho > 0.3$) with attention weights during generation. Validation shows Contriever $\rho = 0.612$ (95\% CI [0.601, 0.624]).

\textbf{H2 (Diversity for Multi-Hop Reasoning)}: Multi-hop QA requires bridging distant facts across semantically dissimilar passages. Diversity-aware eviction preserves contrastive evidence that pure relevance ranking would discard.
